\documentclass[11pt]{article}
\usepackage[a4paper,margin=20mm]{geometry}
\usepackage{amsmath,amssymb,amsthm,microtype}
\usepackage[colorlinks=true,linkcolor=blue,urlcolor=blue]{hyperref}
\newtheorem{lemma}{Lemma}
\newtheorem{theorem}[lemma]{Theorem}
\newtheorem{corollary}[lemma]{Corollary}
\newcommand{\HS}{\mathrm{HS}}
\newcommand{\id}{\mathrm{id}}
\title{Exact polynomial circuits for matrix product unitaries}
\author{TNLean contributors}
\date{2 October 2026}
\begin{document}
\maketitle

Let $N\geq1$ and let $U$ be a unitary on a chain of $N$ qubits. Its consecutive cut rank
at $k$ is the smallest $r$ for which
\[
 U=\sum_{j=1}^r A_j\otimes B_j,
\]
where $A_j$ acts on sites $1,\ldots,k$ and $B_j$ acts on the remaining
sites. An open-boundary MPO of maximal bond dimension two has cut ranks
at most two. The converse follows from successive Schmidt decompositions
of the operator regarded as a vector in Hilbert--Schmidt space.

Styliaris--Trivedi--Cirac, arXiv:2508.08160, give a circuit construction
whose cost depends on a conditioning parameter. Section~5 below removes
this dependence: every MPU of bounded consecutive cut rank has an exact
polynomial-size circuit with clean auxiliary registers. Its proof chooses
positive boundary metrics by determinant maximization over affine spaces
of cap Grams. The first three sections give a separate rank-two argument
with fewer auxiliaries. Section~\ref{sec:uniform-averaged-grams} gives a
preliminary all-length uniform argument by averaging cap Grams.
The relevant source statements are Definition~1 (uniform-bulk MPU),
Assumption~1 and Theorem~1, and the paragraph on nonuniform MPUs ending
with Theorem~2, all in the main text of arXiv:2508.08160v2.
No claim of priority is intended.
Below, ``source lines'' are line numbers of the file \texttt{main.tex} in
the arXiv source archive of version~2,
\url{https://arxiv.org/e-print/2508.08160v2}, and names such as
\nolinkurl{eq:def_q_k} are the equation and lemma labels of that file.

Circuits use arbitrary one- and two-qudit gates. Auxiliary registers begin
and end in $|0\rangle$. Gate bounds are nonuniform existence statements;
they do not assert an efficient classical procedure for obtaining exact
gate parameters from inexact tensor data. The general exact circuit
theorem of Section~5 and its operator-norm approximation corollaries,
including neighboring-pair circuits, exact auxiliary erasure, and the
periodic-boundary specialization, are formalized in Lean in
\nolinkurl{TNLean/MPS/MPU/ExactBoundedCutRankCircuit.lean} and
\nolinkurl{TNLean/MPS/MPU/ApproximateBoundedCutRankCircuit.lean}. The separate rank-two structural argument in the
first three sections is not yet fully formalized.

\section{A two-operator simultaneous singular-value decomposition}

\begin{lemma}\label{lem:two-kraus}
Suppose two square complex matrices $E,F$ satisfy
\[
 E^\dagger E+F^\dagger F=I,
 \qquad EE^\dagger+FF^\dagger=I.
\]
There are unitaries $L,R$ and diagonal matrices $D_E,D_F$ such that
$E=LD_ER$ and $F=LD_FR$.
\end{lemma}
\begin{proof}
Put $H=E^\dagger E$. Then $0\leq H\leq I$ and
$F^\dagger F=I-H$. In finite dimension the partial isometries in the
polar decompositions of square matrices extend to unitaries. Thus
\[
 E=V\sqrt H,\qquad F=W\sqrt{I-H}
\]
for unitaries $V,W$. The second assumed identity gives
\[
 VHV^\dagger+W(I-H)W^\dagger=I,
 \qquad VHV^\dagger=WHW^\dagger.
\]
Consequently the unitary $Q=V^\dagger W$ commutes with $H$.
A commuting Hermitian matrix and unitary matrix are simultaneously
unitarily diagonalizable: diagonalize $H$, then diagonalize the
restriction of $Q$ to each eigenspace of $H$. Write
$H=S\operatorname{diag}(h_j)S^\dagger$ and
$Q=S\operatorname{diag}(q_j)S^\dagger$. Then
\[
 E=VS\operatorname{diag}(\sqrt{h_j})S^\dagger,
 \qquad
 F=VS\operatorname{diag}(q_j\sqrt{1-h_j})S^\dagger.
\]
Taking $L=VS$ and $R=S^\dagger$ proves the assertion, including zero
singular values.
\end{proof}

Unitary polar completion, equality of the two conjugations of $H$, and
the relative-polar commutation conclusion are proved in the Lean module
\nolinkurl{TNLean/MPS/MPU/TwoKrausPolar.lean}. Simultaneous
diagonalization and the full conclusion of Lemma~\ref{lem:two-kraus}
are not yet asserted as Lean theorems.

\begin{lemma}\label{lem:bipartite-diagonal}
Every bipartite unitary of operator Schmidt rank two admits a
factorization
\[
 U=(L_A\otimes L_B)\operatorname{diag}(M_{ij})(R_A\otimes R_B),
\]
where the four outside factors are unitary, all $|M_{ij}|=1$, and
$\operatorname{rank}M=2$. Here $\operatorname{diag}(M_{ij})$ means
the diagonal operator on the product basis indexed by $(i,j)$.
\end{lemma}
\begin{proof}
Let the two dimensions be $d_A,d_B$. Choose a Schmidt expansion
$U=A_1\otimes B_1+A_2\otimes B_2$ such that the $A_i$ are
Hilbert--Schmidt orthogonal and
\[
 \operatorname{Tr}(B_i^\dagger B_j)=d_B\delta_{ij}.
\]
Tracing $U^\dagger U=I$ and $UU^\dagger=I$ over the second factor gives
\[
 A_1^\dagger A_1+A_2^\dagger A_2=I,
 \qquad A_1A_1^\dagger+A_2A_2^\dagger=I.
\]
Lemma~\ref{lem:two-kraus} simultaneously diagonalizes the two $A_i$
by left and right unitaries $L_A,R_A$.

For completeness, the opposite side can be diagonalized without changing
these choices. Put
\[
 \alpha_i=\frac{\|A_i\|_\HS}{\sqrt{d_A}}>0,
 \qquad \widehat A_i=A_i/\alpha_i,
 \qquad \widehat B_i=\alpha_iB_i.
\]
The $\widehat A_i$ are orthogonal and have squared norm $d_A$.
Partial traces over the first factor now show that $\widehat B_1,
\widehat B_2$ satisfy the hypotheses of Lemma~\ref{lem:two-kraus}.
Their simultaneous left and right diagonalization also diagonalizes the
original $B_i$, since each rescaling is scalar. Hence the transformed
$U$ is diagonal in a product basis. Its diagonal entries have modulus
one by unitarity. They form a matrix $M$ of rank two because invertible
local changes preserve operator Schmidt rank, and the Schmidt rank of
this diagonal operator is exactly the ordinary matrix rank of $M$.
\end{proof}

\section{Two sectors with unitary factors in the Schmidt spaces}

\begin{lemma}\label{lem:circle}
If a nonempty finite matrix $M$ has unit-modulus entries and rank at
most two, either its rows have at most two types up to a scalar multiple,
or its columns have at most two such types.
\end{lemma}
\begin{proof}
Rank one is immediate. Otherwise choose independent rows $u,v$.
Divide column $j$ by $u_j$, and put $t_j=v_j/u_j$. Each resulting row
is $a_i+b_it_j$, with $|t_j|=1$. If there are at most two values of
$t_j$, there are at most two projective column types. Otherwise the
identity
\[
 1=|a_i+b_it_j|^2
   =|a_i|^2+|b_i|^2+
      2\operatorname{Re}(a_i\overline{b_i}\overline{t_j})
\]
holds at three distinct points of the unit circle. A nonconstant real
linear functional has a level line, which intersects the unit circle
in at most two points. Thus $a_i\overline{b_i}=0$, and every row is
a scalar multiple of either $1$ or $t$.
\end{proof}

The preceding dichotomy is a mathematical proof. Its Lean formalization
is not claimed in this note.

\begin{lemma}\label{lem:two-sectors}
Let $U$ be a bipartite unitary of Schmidt rank two. After possibly
interchanging the names of the two factors, there are unitaries
$A_0,A_1$ on the first factor and matrices $B_0,B_1$ on the second such
that
\[
 U=A_0\otimes B_0+A_1\otimes B_1.
\]
The matrices $B_+=B_0+B_1$ and $B_-=B_0-B_1$ are unitary, and
\[
 R=B_+^\dagger B_-=R^\dagger,\qquad R^2=I.
\]
Moreover $A_0,A_1$ belong to the first Schmidt space of $U$, and
$B_+,B_-$ belong to its second Schmidt space.
\end{lemma}
\begin{proof}
Apply Lemmas~\ref{lem:bipartite-diagonal} and~\ref{lem:circle}.
Suppose, for definiteness, that there are two projective row types.
Choose normalized representative phase rows $r_0,r_1$ and write
$M_{ij}=g_i r_{p(i),j}$, with $|g_i|=|r_{s,j}|=1$ and
$p(i)\in\{0,1\}$. Both classes are nonempty. Let $P_s$ be the
diagonal projection onto $p(i)=s$. The original $U$ is
\[
 \sum_{s=0}^1
 \left[L_A\operatorname{diag}(g_i)P_sR_A\right]\otimes
 \left[L_B\operatorname{diag}(r_{s,j})R_B\right].
\]
Rename the second factor as the first, so that the two bracketed phase
operators there are $A_0,A_1$, while the former first-factor operators
are $B_0,B_1$. Then
\[
 B_+=L_A\operatorname{diag}(g_i)R_A,
 \qquad
 B_-=L_A\operatorname{diag}(g_i)(P_0-P_1)R_A.
\]
Both are unitary, and
$R=R_A^\dagger(P_0-P_1)R_A$ is explicitly a Hermitian involution.
The column case is identical with the roles reversed. Each displayed
pair is independent because $U$ has Schmidt rank two. Thus the pairs
span the original Schmidt spaces, proving the final assertion.

In particular the Hermitian property of $R$ follows from this sector
construction. It does not follow from the unitarity of $B_+,B_-$ alone.
\end{proof}

\begin{lemma}\label{lem:inheritance}
Let an operator $U$ have rank at most two at every consecutive cut of
a chain. Split the chain into two consecutive intervals. Every element
of either Schmidt space at this split has rank at most two at every
consecutive cut inside its interval.
\end{lemma}
\begin{proof}
Every element of a Schmidt space is obtained by applying a linear
functional on the opposite factor to $U$. At a cut internal to the
retained interval, apply that same functional to a rank-two expansion
of $U$. It still has at most two summands. No assertion of positivity or
physical realizability of the functional is needed; the statement is
solely a bound on operator Schmidt rank.
\end{proof}

Taking adjoints preserves operator Schmidt rank: the adjoints of the
factors in a rank-$r$ expansion give a rank-$r$ expansion of the adjoint,
and applying the same argument a second time gives equality. This also
preserves every consecutive cut-rank bound.

\section{A recursion with a fixed number of controls}

Write $C_c(W)$ for $W$ applied when a control qubit $c$ equals one.
For the decomposition in Lemma~\ref{lem:two-sectors}, put
$P_s=(I+(-1)^sR)/2$. Then $B_s=B_+P_s$, and
\[
 U=(I\otimes B_+)(A_0\otimes P_0+A_1\otimes P_1).
\]
Introduce a clean qubit $z$. The unitary
\[
 E_c=H_z C_{cz}(R)H_z
\]
records the sector in $z$ when $c=1$, and is the identity when $c=0$.
Indeed, for a vector in the image of $P_s$ and $z$ initially zero,
it sets $z=s$ if $c=1$. Since $R$ is a Hermitian involution,
$E_c^\dagger=E_c$.

After $E_c$, apply $A_0$ controlled by $c(1-z)$ and $A_1$ controlled
by $cz$. These operations act on the opposite interval and preserve
the recorded sector. Erasing the record and applying $B_+$ controlled
by $c$ therefore implements $C_c(U)$ with $z$ returned to zero.
This already uses seven controlled child unitaries: two for recording,
two for erasing, two for the $A_s$, and one for $B_+$.

One child call can be saved. Since $R=B_+^\dagger B_-$, and the
outer control commutes with $H_z$,
\[
 C_c(B_+)E_c^\dagger
   =H_z C_{c(1-z)}(B_+)C_{cz}(B_-)H_z.
\]
To verify this identity, use $E_c^\dagger=E_c$ and cancel the adjacent
$C_c(B_+)$ and $C_{cz}(B_+^\dagger)$. On $c=0$ all controlled factors
are identities; on $c=1,z=0$ the product is $B_+$, and on $c=1,z=1$
it is $B_-$. Hence a controlled implementation uses precisely the
following six child calls:
\[
 \underbrace{B_-,\ B_+^\dagger}_{\text{sector recording}},
 \qquad A_0,\ A_1,
 \qquad
 \underbrace{B_+,\ B_-}_{\text{sector erasure and final factor}}.
\]
The last two calls lie between Hadamards on $z$, as in the displayed
identity. Each child call has one control. Conjunctions $cz$ and
$c(1-z)$ are computed into two additional clean qubits by Toffoli
gates, used as the controls, then erased before any subsequent
Hadamard on $z$. Toffoli gates have constant-size exact decompositions
into the allowed gates. Thus three auxiliary qubits at this level
suffice, and the number of controls does not grow with recursion depth.
The control may be entangled with the data or with other registers;
the identities are operator identities and remain valid in that case.

The clean-record merging identity and the cancellation combining sector
erasure with the final unitary factor are proved in the Lean module
\nolinkurl{TNLean/MPS/MPU/TwoSectorMerging.lean}. Its hypotheses
explicitly supply an orthogonal projection or a Hermitian involution.
These algebraic identities support the six-call construction; the
structural decomposition of an arbitrary Schmidt-rank-two unitary and
its rank-two recursive circuit bound remain unformalized mathematical proofs.

\begin{theorem}\label{thm:circuits}
For $N\geq1$, every $N$-qubit unitary of rank at most two at every consecutive cut
has an exact circuit with
\[
 O(N^{\log_2 6})\quad\text{one- and two-qubit gates},
 \qquad O(\log(N+1))\quad\text{clean auxiliary qubits}.
\]
A sequential implementation has depth $O(N^{\log_2 6})$.
On a line, routing gives gate count and depth
$O(N^{1+\log_2 6})$, with the same auxiliary-space bound.
\end{theorem}
\begin{proof}
Construct the controlled unitary recursively. At a balanced split of
rank two, use the six-call construction just proved. Its four child
unitaries $A_0,A_1,B_+,B_-$ have consecutive ranks at most two by
Lemma~\ref{lem:inheritance}. In particular, the involution $R$ need
not be assigned its own recursive implementation: its two factors
$B_+^\dagger,B_-$ are the child unitaries that are implemented.
The adjoint $B_+^\dagger$ retains the same internal cut ranks; its
controlled implementation can equivalently be obtained by reversing the
controlled circuit for $B_+$.

At a rank-one split, $U=A\otimes B$ with $A,B$ rescaled to be
unitary. This rescaling is possible because
$A^\dagger A\otimes B^\dagger B=I$ forces both positive factors to
be scalar multiples of the identity. The two controlled child calls
also obey Lemma~\ref{lem:inheritance}. At one site the controlled
unitary is an allowed two-qubit gate.

If $G(n)$ denotes the maximum gate count for such controlled circuits,
the balanced recursion gives
\[
 G(n)\leq6G(\lceil n/2\rceil)+C
 \quad(n\geq2),
\]
with an absolute constant $C$ for the Hadamards and guard computations.
Its recursion tree has height at most $\lceil\log_2 n\rceil$ and
$O(6^{\lceil\log_2 n\rceil})=O(n^{\log_2 6})$ nodes. This height is not
a circuit depth: the six child calls at each node run one after another
on shared auxiliary workspace, so the circuit depth obeys the same
recurrence as the gate count and is bounded by it, not by
$O(\log n)$.
The three current-level auxiliary qubits are retained during a child
call, while the child's space is reused for subsequent calls. Thus at
most $3\lceil\log_2 n\rceil+O(1)$ auxiliary qubits are present.
Initialize the root control to one and return it to zero at the end.

Sequential depth is bounded by the gate count. There are
$N+O(\log(N+1))=O(N)$ total qubits. On a line every two-qubit
interaction can be routed using $O(N)$ adjacent swaps and then the
positions restored, giving the stated gate and depth bounds. No copying
of quantum data is used or required.
\end{proof}

The numerical estimate for a recursion with a fixed number of child
calls is formalized separately by the Lean theorem
\nolinkurl{MPUCircuit.cost_le_pow_of_constant_branching} in
\nolinkurl{TNLean/MPS/MPU/ConstantBranchingBounds.lean}. Its hypotheses
concern a numerical recurrence; it does not assert that the controlled
unitary decomposition above has been formalized. The exact finite-matrix
identity for the six-call merging is checked in
\nolinkurl{scripts/check_mpu_rank_two_controlled_merging.py}, on several
noncommuting examples and both orientations of the controlling interval.

The potentially complicated local unitaries in
Lemma~\ref{lem:bipartite-diagonal} are never implemented as circuits.
They establish the existence of four unitary elements of the original
Schmidt spaces. The recursion implements those elements directly, using
their inherited rank bounds. This distinction prevents a circular
assumption that the local diagonalizing factors are already efficient.

The exact theorem also implies the operator-norm approximation statement
for this rank-two class, since the exact circuit has zero error. It
does not require any conditioning hypothesis. A preliminary uniform
argument follows; Section~5 then establishes the general fixed-bond
exact theorem by a different choice of positive metrics.

\section{Uniform families and averaged boundary Grams}
\label{sec:uniform-averaged-grams}

The all-length condition in the definition of a uniform-bulk MPU
allows boundary Grams from different cap lengths to be averaged.
Their isometry identities are separately linear in the two Grams.
This gives a familywise polynomial bound without the source's
exact-length spanning assumption.

\begin{theorem}\label{thm:uniform-averaged-grams}
Fix a physical dimension $d\geq2$, a finite-dimensional space $E$,
matrices $A^{ij}\in\operatorname{End}(E)$, a row $l\in E^*$, and
a column $r\in E$, where $1\leq i,j\leq d$. Suppose that
\[
 (U_N)_{\mathbf i,\mathbf j}
   =lA^{i_1j_1}\cdots A^{i_Nj_N}r
\]
is unitary for every $N\geq1$. Then $U_N$ has an exact circuit of
polynomial size and depth, using arbitrary one- and two-qudit gates
and $O(N)$ auxiliary qudits which begin and end in $|0\rangle$.
The polynomial constants and exponent may depend on the fixed
$A,l,r,d$. The same conclusion holds with gates between adjacent sites
on a line. No exact-length full-rank condition on the cap Grams
is required.
\end{theorem}

This is the fixed-presentation, all-length hypothesis of the
Uniform-bulk MPU definition in Styliaris--Trivedi--Cirac,
arXiv:2508.08160. The circuit step uses their Local isometries:
Uniform case lemma and Merging lemma, in the Supplemental Material
of arXiv:2508.08160v2. The additional step
here is to take convex averages of cap Grams over different lengths
before taking their square roots.

\begin{proof}
\subsection*{The positive-word representation}

Write $\alpha=(i,j)$ for the physical pair alphabet and
$A_w=A_{\alpha_1}\cdots A_{\alpha_m}$ for a word $w$ of length $m$.
Put
\[
 E_+=\operatorname{span}\{A_wr:|w|\geq1\},\qquad
 K=\{x\in E_+:lA_wx=0\text{ for every }|w|\geq1\},
 \qquad W=E_+/K.
\]
Both $E_+$ and $K$ are invariant under every $A_\alpha$.
For $K$, this follows from
$lA_wA_\alpha x=lA_{w\alpha}x=0$ for every positive-length $w$.
Thus $A_\alpha$ induces a matrix $\overline A_\alpha$ on $W$.
Define the endpoint tensors
\[
 B_\alpha([x])=lA_\alpha x,
 \qquad C_\alpha=[A_\alpha r].
\]
The row $B_\alpha$ is well defined because $lA_\alpha$ vanishes on
$K$, and $A_\alpha r$ lies in $E_+$. For every $N\geq2$,
\[
 (U_N)_{\mathbf i,\mathbf j}
 =B_{(i_1,j_1)}\overline A_{(i_2,j_2)}\cdots
       \overline A_{(i_{N-1},j_{N-1})}C_{(i_N,j_N)}.
\]
The case $N=2$ means $B_{(i_1,j_1)}C_{(i_2,j_2)}$.
No assertion that $l$ or $r$ themselves descend to $W$ is needed.
The one-site unitary will be treated separately.

For every positive word define a row and column on $W$ by
\[
 f_u([x])=lA_ux,\qquad q_v=[A_vr].
\]
The columns $q_v$ span $W$ by construction, and the rows $f_u$
span $W^*$: their common kernel is zero by the definition of $K$.
Let $w=\dim W\leq\dim E$. Here $w>0$, since otherwise $U_2=0$.

In fact cumulative word lengths $1,\ldots,w$ suffice on each side.
The cumulative column spaces $S_j=\operatorname{span}\{q_v:1\leq|v|\leq j\}$
satisfy
\[
 S_{j+1}=S_1+\sum_\alpha\overline A_\alpha S_j.
\]
If $S_{j+1}=S_j$, then $S_j$ is invariant and contains every
future reachable column, so $S_j=W$. Since $S_1\ne0$, dimension
therefore increases strictly until it reaches $W$, in at most
$w$ steps. The same argument for rows, using right multiplication
by $\overline A_\alpha$, gives
\[
 \operatorname{span}\{f_u:1\leq|u|\leq w\}=W^*.
\]

\subsection*{Averaging the cap Grams}

Choose a Hilbert-space structure on $W$ and set, for $m,n\geq1$,
\[
 P_m=d^{-m}\sum_{|u|=m}f_u^\dagger f_u,
 \qquad
 Q_n=d^{-n}\sum_{|v|=n}q_vq_v^\dagger.
\]
The sums run over all words in the pair alphabet: the factors
$d^{-m},d^{-n}$ average the cap input indices and sum the cap
output indices. Define
\[
 P=\frac1w\sum_{m=1}^wP_m,
 \qquad Q=\frac1w\sum_{n=1}^wQ_n.
\]
The cumulative spanning statements imply $P,Q\succ0$.
Indeed a vector in the kernel of $P$ is annihilated by every
$f_u$ of length at most $w$, and a vector in the kernel of $Q$
is orthogonal to every $q_v$ of those lengths.

For a central interval of length $k\geq1$, put
\[
 T_{\mathbf a,\mathbf b}
 =\overline A^{a_1b_1}\cdots\overline A^{a_kb_k},
 \qquad
 \Phi_k(P',Q')_{\mathbf b,\mathbf b'}
 =\sum_{\mathbf a}\operatorname{Tr}\!\left(
 T_{\mathbf a,\mathbf b}^\dagger P'
 T_{\mathbf a,\mathbf b'}Q'\right).
\]
For each $m,n\geq1$, unitarity at length $m+k+n$ gives
\[
 \Phi_k(P_m,Q_n)_{\mathbf b,\mathbf b'}
       =\delta_{\mathbf b,\mathbf b'}.
\]
To see the normalization explicitly, fix the input words on the
two caps, sum over both cap outputs and the central output in
$U_{m+k+n}^\dagger U_{m+k+n}=\id$, and average over the cap inputs
with weights $d^{-m}$ and $d^{-n}$. The resulting left side is
\[
 d^{-m-n}\sum_{\mathbf a,u,v}
 \overline{f_uT_{\mathbf a,\mathbf b}q_v}\,
              f_uT_{\mathbf a,\mathbf b'}q_v,
\]
which is the displayed trace formula.
The function $\Phi_k$ is separately linear in $P',Q'$.
Consequently the independent averages over $m$ and $n$ give
\[
 \Phi_k(P,Q)=\id\qquad(k\geq1).
\]
It is the Grams which are averaged here, not their square roots.

The boundary versions follow in the same way. Write
\[
 F_{\mathbf a,\mathbf b}([x])
   =lA^{a_1b_1}\cdots A^{a_kb_k}x,
 \qquad
 G_{\mathbf a,\mathbf b}
 =[A^{a_1b_1}\cdots A^{a_kb_k}r].
\]
Tracing only a right cap in $U_{k+n}^\dagger U_{k+n}=\id$ gives
\[
 \sum_{\mathbf a}\operatorname{Tr}\!\left(
 F_{\mathbf a,\mathbf b}^\dagger F_{\mathbf a,\mathbf b'}Q_n\right)
 =\delta_{\mathbf b,\mathbf b'},
\]
and tracing only a left cap in $U_{m+k}^\dagger U_{m+k}=\id$ gives
\[
 \sum_{\mathbf a}G_{\mathbf a,\mathbf b}^\dagger
                   P_mG_{\mathbf a,\mathbf b'}
 =\delta_{\mathbf b,\mathbf b'}.
\]
The same identities hold after averaging to $Q$ and $P$.

\subsection*{The interval isometries and their merger}

Take $L=P^{1/2}$ and $R=Q^{1/2}$. Regard the row and column indices
of $LT_{\mathbf a,\mathbf b}R$ as two additional output registers.
The maps with amplitudes
\[
 LT_{\mathbf a,\mathbf b}R,\qquad
 F_{\mathbf a,\mathbf b}R,\qquad
 LG_{\mathbf a,\mathbf b}
\]
are respectively the interior, left-boundary and right-boundary
interval isometries. For the first, the Gram matrix is
\[
 \sum_{\mathbf a}\operatorname{Tr}\!\left(
 (LT_{\mathbf a,\mathbf b}R)^\dagger
  LT_{\mathbf a,\mathbf b'}R\right)
 =\Phi_k(P,Q)_{\mathbf b,\mathbf b'},
\]
and the preceding one-boundary identities prove the other two cases.
For $k=1$ all these maps have fixed dimension, so each has a unitary
extension with a fixed number of one- and two-qudit gates and
auxiliary qudits. Non-power-of-$d$ flag dimensions are embedded in
a fixed number of $d$-dimensional registers.

For adjacent intervals contract the right flag of the first with
the left flag of the second using the coefficient matrix
\[
 T=R^{-1}L^{-1},\qquad
 \mu=\|T\|_{\HS}<\infty.
\]
Then $RTL=\id$, so this contraction gives precisely the longer
interval tensor. The corresponding operator resets the two
contracted registers to zero. Its norm is $\mu$, since its only
nonzero output row is the contraction bra. In these matrix
conventions $\mu^2=\operatorname{Tr}(P^{-1}Q^{-1})$;
defining $\mu$ by the actual contraction bra avoids any dependence
on vectorization transpose conventions.

The product of the two child maps is an isometry and the merged
map is an isometry. Hence the contraction has norm one on the
child image, and in particular $\mu\geq1$. Its fixed-dimensional
unitary dilation implements the contraction divided by $\mu$.
On that image the success probability is $\mu^{-2}$, independent
of the unknown physical input. The source's Subspace amplitude
amplification lemma and Merging lemma therefore apply.

For completeness, the exact amplification can be chosen using
only ordinary plane rotations. Put
\[
 \ell=\left\lceil\frac{\pi\mu}{4}\right\rceil,
 \qquad \theta=\frac{\pi}{4\ell+2},
 \qquad c=\mu\sin\theta\leq1.
\]
Add an attenuation register initialized to zero and rotate it so
that its zero amplitude is $c$. Let $\mathcal D$ be the merger
dilation together with this rotation. On each vector in the
initial child image it now has the form
\[
 \mathcal D\psi
   =\sin\theta\,g_\psi+\cos\theta\,h_\psi,
\]
where $g_\psi$ is the desired merged output with the contracted
flags and both new registers zero, and $h_\psi$ is orthogonal to
the corresponding success subspace. Both $g_\psi$ and $h_\psi$
preserve inner products as $\psi$ ranges over the child image:
this follows from the constant success amplitude and unitarity.

Let $R_{\rm initial}$ reflect about the initial child image with
the new registers zero, and let $R_{\rm success}$ reflect about
the zero contracted and new registers. The operator
\[
 -\mathcal D R_{\rm initial}\mathcal D^\dagger R_{\rm success}
\]
rotates each plane spanned by $g_\psi,h_\psi$ through $2\theta$.
After $\ell$ rotations the output is exactly $g_\psi$, since
$\sin((2\ell+1)\theta)=1$. The reflection about the initial image
is obtained by undoing the child unitary extensions, reflecting
about their initialized auxiliaries and the new zero registers,
and reapplying the extensions. Thus a bounded number of child
circuit invocations, depending only on $\mu$, performs the merge.
The contracted registers and the amplification auxiliaries return
exactly to zero; the interval's outer flags are retained.

\subsection*{Circuit bounds and exact erasure}

A balanced tree merges the single-site isometries into $U_N$.
If $H(n)$ denotes the gate cost of an interval of length $n$,
the preceding implementation gives a recurrence
\[
 H(n)\leq c_\mu\bigl(H(\lceil n/2\rceil)+H(\lfloor n/2\rfloor)\bigr)
            +c_\mu n,
\]
where $c_\mu$ also includes the fixed physical and flag dimensions.
The linear term covers the image reflections and their reversible
conjunctions. The recurrence is polynomial; for example any fixed
$p>1+\log_2c_\mu$ gives $H(n)=O(n^p)$ after increasing the constant.
Depth is at most the gate count and is consequently polynomial.

Each leaf and each merge introduces a fixed number of logical
auxiliary registers, so the tree uses $O(N)$ of them. The reversible
conjunction scratch can be taken from a common $O(N)$-sized pool:
each conjunction is erased before a child invocation, and each
child returns this scratch to zero for every logical input. Thus
amplification repetitions and nested calls reuse the same pool.
At every internal merge the two inner flag registers are reset
exactly; only the interval's external flags remain. At the root
the boundary tensors have no external flags. The final map is
therefore $U_N$ with every auxiliary register exactly zero.
For $N=1$, apply the fixed one-qudit unitary $U_1$ directly.
On a line, route each two-qudit interaction by adjacent swaps and
restore the positions. There are $O(N)$ total qudits, so this adds
at most an $O(N)$ factor to the gate and depth bounds.
\end{proof}

The averaging proof has familywise quantifiers: its bound depends
on the fixed presentation through $P,Q$ and $\mu$. Section~5 gives
a stronger result by optimizing over an affine space of positive
metrics, and covers nonuniform tensors specified at a single length.
The averaging argument above is a mathematical proof; this separate
familywise circuit construction is not asserted to have been formalized
in Lean. The stronger general construction of Section~5 is formalized
in Lean.

Fixed trace boundaries are included. If the coefficient is
$\operatorname{Tr}(bA_{\alpha_1}\cdots A_{\alpha_N})$, take the enlarged
bond space $\operatorname{End}(E)$, let each letter act by
$X\mapsto A_\alpha X$, and take the right vector $b$ and left
functional $\operatorname{Tr}$. Cyclicity of the trace gives the
same coefficients. This includes periodic boundaries $b=I$ and
requires no invertibility assumption on $b$.

An exact-length loss of rank does not obstruct the averages.
For example, take a fixed unit-modulus phase $\lambda$ of infinite order and the
one-qubit projection $P_0=|0\rangle\langle0|$. The family
\[
 U_N=\lambda^N\id+(1-\lambda^N)P_0^{\otimes N}
\]
is unitary for every $N$. It has a fixed three-dimensional
presentation
\[
 A^{ij}=\operatorname{diag}(\lambda\delta_{ij},
                  (P_0)_{ij},\lambda(P_0)_{ij}),\qquad
 l=(1,1,-1),\quad r=(1,1,1)^T.
\]
At each fixed positive word length the reachable columns span
at most two dimensions: a nonzero word containing the pair
$(1,1)$ gives a multiple of $(1,0,0)^T$, while the all-$(0,0)$
word gives
$(\lambda^m,1,\lambda^m)^T$. The corresponding rows also span
at most two dimensions. But lengths one and two together span
all three dimensions on both sides, so their averaged Grams are
positive definite. This loss of rank occurs in a minimal fixed
presentation: with $a=(0,0)$ and $b=(1,1)$, the matrix
$(lA_uA_vr)_{u,v\in\{a,aa,b\}}$ has determinant
$-\lambda^4(\lambda-1)^2\ne0$. Every nontrivial consecutive
cut of $U_N$ has rank two for $N\geq2$, since
$1-\lambda^N\ne0$. Thus the deficient exact-length cap ranks
persist on nondegenerate data; they are repaired by averaging, rather
than by assuming exact-length full rank. The proof uses unitarity at
every cap-extended length.

\clearpage
\section{Determinant normalization for arbitrary finite MPUs}
\label{sec:determinant-normalization}

The interval-isometry construction can be normalized without a lower
bound on the nonzero operator Schmidt values. The relevant freedom is
larger than the convex set of Grams obtained from physical cap states:
one may choose a positive matrix in their real affine hull. Maximizing
its determinant gives a dual normalization at the same cut. These
normalizations are compatible across intervals, including intervals
with entangled physical input states.

\begin{theorem}\label{thm:general-determinant-circuits}
Fix a physical dimension $d\geq2$ and an integer $D\geq1$.
Let $U$ be a unitary on $N\geq1$
$d$-dimensional sites, with operator Schmidt rank at most $D$ across
each consecutive cut. There is an exact circuit for $U$, using
arbitrary one- and two-qudit gates, with
\[
 (D+1)^{O_d(1)}N^{O_d(\log(D+1))}
\]
gates and $O_d(N\log(D+1))$ auxiliary qudits. Every auxiliary begins
and ends in $|0\rangle$. The same bounds, with a change in the numerical
exponent, hold using only adjacent two-qudit gates on a line.
The implementation is an operator equality on clean-auxiliary inputs,
and consequently remains exact when the data are entangled with an
arbitrary reference system.
\end{theorem}

The theorem addresses the nonuniform fixed-bond-dimension circuit
question treated under a spectral or conditioning requirement in
Styliaris--Trivedi--Cirac, arXiv:2508.08160v2, in particular the
definition of $q_k$ and Theorem~2 in the main-text paragraph on
nonuniform MPUs (source lines 1373--1402).
The proof uses the isometry and merging argument of that paper,
but obtains the positive metrics from affine Gram spaces rather than
from physical cap states. The main-text definition of $q_k$ ranges over
valid isometry normalizers (lines 1375--1379), whereas its appendix
restatement restricts the boundary squares to physical cap-Gram sets
(lines 2123--2124). These are different domains. We compute the merger
norm directly below, without identifying it with either infimum. Any
comparison must specify the domain and match the paper's transpose
and boundary-index conventions.

\begin{proof}
\subsection*{A minimal open-boundary representation}

Successive Schmidt decompositions of the operator give an open MPO
\[
 U_{\mathbf a,\mathbf b}
 =A_1^{a_1b_1}\cdots A_N^{a_Nb_N},
 \qquad A_j^{ab}:E_j\longrightarrow E_{j-1},
\]
where $E_0=E_N=\mathbb C$ and
$r_j=\dim E_j\leq D$ for $1\leq j<N$.
At every cut the prefix rows and suffix columns span $E_j^*$ and
$E_j$, respectively. This follows from minimality of the cut rank:
a linear dependence on either side would reduce that rank.
All $r_j$ are positive, and put $r_0=r_N=1$.

At cut $j$, denote the prefix row by
\[
 f^{(j)}_{\mathbf a,\mathbf b}
   =A_1^{a_1b_1}\cdots A_j^{a_jb_j}.
\]
For a density matrix $\rho$ on the first $j$ sites, define its
left cap Gram by
\[
 P_j(\rho)=
 \sum_{\mathbf a,\mathbf b,\mathbf b'}
 \rho_{\mathbf b',\mathbf b}
 (f^{(j)}_{\mathbf a,\mathbf b})^\dagger
  f^{(j)}_{\mathbf a,\mathbf b'}.
\]
This is positive semidefinite. For a pure input $\rho=|\psi\rangle
\langle\psi|$ it is the sum of the positive matrices
$(\sum_{\mathbf b}\psi_{\mathbf b}f^{(j)}_{\mathbf a,\mathbf b})^\dagger
 (\sum_{\mathbf b}\psi_{\mathbf b}f^{(j)}_{\mathbf a,\mathbf b})$.
The map $\rho\mapsto P_j(\rho)$ is linear on Hermitian matrices.

Let $g^{(j)}_{\mathbf c,\mathbf e}$ be the suffix column obtained
from sites $j+1,\ldots,N$, and let
\[
 \overline Q_j=d^{-(N-j)}
       \sum_{\mathbf c,\mathbf e}
        g^{(j)}_{\mathbf c,\mathbf e}
       (g^{(j)}_{\mathbf c,\mathbf e})^\dagger.
\]
The spanning of the suffix columns gives $\overline Q_j\succ0$.
Likewise the maximally mixed prefix input gives
\[
 P_j(d^{-j}\id)\succ0.
\]
Unitarity, applied to the product of $\rho$ with the maximally
mixed suffix input, implies
\[
 \operatorname{Tr}(P_j(\rho)\overline Q_j)=1
       \qquad\text{for every density matrix }\rho.
\]
The endpoint cases are scalar: $P_0(1)=1$, while
$P_N(\rho)=\operatorname{Tr}(U\rho U^\dagger)=1$.

\subsection*{The determinant maximizer and its dual metric}

Let $\mathcal C_j$ be the real affine hull of all $P_j(\rho)$ as
$\rho$ ranges over density matrices on the prefix. This is an
affine subspace of the Hermitian matrices. Its elements are finite
real combinations
\[
 P=\sum_t a_tP_j(\rho_t),\qquad
 a_t\in\mathbb R,\qquad\sum_ta_t=1.
\]
The coefficients need not be positive. Thus $\mathcal C_j$ is
distinguished from the physical convex set of cap Grams.
The preceding trace identity extends to every $P\in\mathcal C_j$.

The set $\mathcal C_j\cap\{P\succeq0\}$ is compact. It is closed,
and $\overline Q_j\succ0$ gives
\[
 1=\operatorname{Tr}(P\overline Q_j)
     \geq\lambda_{\min}(\overline Q_j)\operatorname{Tr}P,
\]
which bounds the positive matrices. It contains the positive
definite maximally mixed cap Gram. A determinant maximizer $P_j^*$
therefore exists and has positive determinant, hence $P_j^*\succ0$.

For every Hermitian direction $H$ parallel to $\mathcal C_j$,
the matrices $P_j^*+tH$ remain positive definite for all sufficiently
small positive and negative $t$. Jacobi's derivative formula gives
\[
 0=\left.\frac{d}{dt}\log\det(P_j^*+tH)\right|_{t=0}
   =\operatorname{Tr}((P_j^*)^{-1}H).
\]
Set
\[
 Q_j^*=\frac1{r_j}(P_j^*)^{-1}.
\]
For every $P\in\mathcal C_j$, stationarity and
$\operatorname{Tr}(P_j^*Q_j^*)=1$ now imply
\[
 \operatorname{Tr}(P Q_j^*)=1.
\]
At the endpoints $\mathcal C_0=\mathcal C_N=\{1\}$, so choose
$P_0^*=Q_0^*=P_N^*=Q_N^*=1$.
Negative affine coefficients have only been used to establish
linear identities. Both final matrices $P_j^*$ and $Q_j^*$ are
positive definite.

\subsection*{Compatibility across every interval}

For $0\leq j<k\leq N$, write
\[
 A_{j:k}^{\mathbf a,\mathbf b}
 =A_{j+1}^{a_{j+1}b_{j+1}}\cdots A_k^{a_kb_k}.
\]
For any density matrix $\sigma$ on this interval, define
\[
 \mathcal T_\sigma(P)=
 \sum_{\mathbf a,\mathbf b,\mathbf b'}
 \sigma_{\mathbf b',\mathbf b}
 (A_{j:k}^{\mathbf a,\mathbf b})^\dagger
 P A_{j:k}^{\mathbf a,\mathbf b'}.
\]
For a prefix density matrix $\rho$, concatenation gives exactly
\[
 \mathcal T_\sigma(P_j(\rho))=P_k(\rho\otimes\sigma).
\]
The input $\sigma$ is arbitrary: it may be entangled among all
sites of the interval. Linearity and the affine-coefficient sum
one therefore give
\[
 \mathcal T_\sigma(\mathcal C_j)\subseteq\mathcal C_k,
 \qquad
 \operatorname{Tr}\bigl(Q_k^*\mathcal T_\sigma(P_j^*)\bigr)=1.
\]

Put
\[
 L_j=(P_j^*)^{1/2},\qquad R_j=(Q_j^*)^{1/2}.
\]
Define $V_{j:k}$ from the physical input on sites $j+1,\ldots,k$
to the physical output and two flag registers by amplitudes
\[
 (V_{j:k})_{(\mathbf a,\alpha,\beta),\mathbf b}
    =(L_j A_{j:k}^{\mathbf a,\mathbf b}R_k)_{\alpha\beta}.
\]
Its input Gram has entries
\[
 (V_{j:k}^\dagger V_{j:k})_{\mathbf b,\mathbf b'}
 =\sum_{\mathbf a}\operatorname{Tr}\!\left(
 (A_{j:k}^{\mathbf a,\mathbf b})^\dagger
 P_j^* A_{j:k}^{\mathbf a,\mathbf b'}Q_k^*\right).
\]
The trace against every density matrix $\sigma$ is one, by the
transfer identity. A Hermitian matrix with that expectation for
every density matrix is the identity. Thus $V_{j:k}$ is an
isometry for every interval. When $j=0$ or $k=N$ the corresponding
flag is one dimensional; $V_{0:N}=U$.

This argument includes both endpoints and arbitrary interval input
states. It uses only the unitarity of the given finite $U$;
unitarity at other lengths is not required.

\subsection*{A Bell contraction of bounded norm}

At a cut $j$, the two adjacent flags have dimension $r_j$.
The matrix that contracts them is
\[
 R_j^{-1}L_j^{-1}
 =\sqrt{r_j}\,(P_j^*)^{1/2}(P_j^*)^{-1/2}
 =\sqrt{r_j}\,\id.
\]
Consequently the merger is
\[
 M_j=\sqrt{r_j}\,|0,0\rangle
                 \sum_{x=0}^{r_j-1}\langle x,x|,
 \qquad \|M_j\|=r_j\leq D.
\]
Applying it to $V_{a:j}\otimes V_{j:b}$ gives exactly $V_{a:b}$
with the contracted flags reset to zero. This follows by inserting
$R_j(R_j^{-1}L_j^{-1})L_j=\id$ between the interval matrices.
The product child map and the target are isometries.

A fixed-dimensional unitary dilation of $M_j/r_j$ has success
amplitude $1/r_j$ on the entire child image. The exact subspace
amplitude amplification used in the preceding section therefore
performs the merge with $O(r_j)$ invocations. Explicitly, take
\[
 \ell_j=\left\lceil\frac{\pi r_j}{4}\right\rceil,
 \qquad\theta_j=\frac{\pi}{4\ell_j+2},
 \qquad c_j=r_j\sin\theta_j\leq1.
\]
An additional attenuation register gives success amplitude
$\sin\theta_j$. Before dilation, reflection about the child image
is implemented by undoing the child unitary extensions, reflecting
about their initialized flags and the new registers, and reapplying
the extensions. Equivalently, let $Z_0$ be the full preparation,
including the children, joining dilation, and attenuation. Reflection
about its prepared image is $Z_0R_0Z_0^\dagger$, where $R_0$ tests
all selected initialized flags for zero. This reflection undoes and
reapplies the small-packet preparation as well as both child circuits.
Reflection about success tests the contracted and new registers for zero. The usual two-plane calculation gives
success amplitude $\sin((2\ell_j+1)\theta_j)=1$ after $\ell_j$
rotations. Thus the merger is exact, including all input-dependent
phases, and all contracted and new registers are erased exactly.
This is the source's Merging lemma applied to the isometries proved
above; no physical-cap representation of the metrics is needed.

\subsection*{Gate and auxiliary bounds}

Encode an $r$-dimensional flag in $\lceil\log_d r\rceil$ qudits.
A single-site isometry acts between a $d$-dimensional input and
a space of dimension at most $d^3(D+1)^2$, after padding the flags.
Complete its columns to a unitary. QR elimination of a unitary of
dimension $m$ uses $O(m^2)$ two-level factors and diagonal phases.
A path between two basis words reduces each two-level factor to
a pattern-controlled one-qudit gate and basis-word transpositions.
These have polynomial cost in the number of qudits: compute the
control conjunction with clean work registers, apply the controlled
gate, and erase the conjunction. Toffoli on two encoded Boolean
controls has a constant-size two-qudit decomposition. All phases
of the two-level factors and the final diagonal are retained.
Thus the leaf extensions and the fixed-dimensional merger
dilations have $(D+1)^{O_d(1)}$ gate cost.

The work registers of these elementary routines return to zero
for every state of the data and logical flags, including arbitrary
nonzero flags. Each implemented extension is therefore a genuine
unitary on the logical registers during amplification and during
inverse child calls. Reflection about zero on $a$ logical flags
uses $O_d(a)$ gates and clean work registers, by a reversible
conjunction; its work also returns to zero for every logical input.

A balanced tree has height $h=\lceil\log_2N\rceil$. The merger
bound gives
\[
 H(n)\leq b_D\bigl(H(\lceil n/2\rceil)+H(\lfloor n/2\rfloor)\bigr)
           +(D+1)^{O_d(1)}n\log(D+1),
 \qquad b_D\leq C_d(D+1).
\]
Increasing $b_D$ to at least two, expansion by tree levels yields
\[
 H(N)\leq(D+1)^{O_d(1)}N b_D^h
          =(D+1)^{O_d(1)}N^{O_d(\log(D+1))}.
\]
Each leaf and each merge contributes $O_d(\log(D+1))$ logical
auxiliary qudits, giving $O_d(N\log(D+1))$ in total. Use a common
work pool of this same order for the conjunctions and gate
decompositions. Work is erased before each child invocation, so
repeated amplification and nested calls reuse the pool. Separate
persistent work pools at every tree level are unnecessary.
At the root there are no external flags, and every other register
has been reset. The circuit therefore implements $U$ with exactly
clean auxiliary output. Since these are equal linear maps, the
statement remains true after tensoring with any reference identity.

To obtain the line version, arrange the data and auxiliary registers
on a line. Route each nonadjacent two-qudit gate using swaps, apply
it, and undo the swaps. There are $O_d(N\log(D+1))$ registers, so
this multiplies the gate count by at most that factor. The auxiliary
count is unchanged, and the exponent remains $O_d(\log(D+1))$.
Depth is bounded by gate count. For $N=1$, apply $U$ directly.
\end{proof}

A periodic MPO of bond dimension $\chi$ can be opened with bond
dimension at most $\chi^2$. Applying the theorem with $D=\chi^2$
gives the same fixed-bond polynomial conclusion, with gate exponent
$O_d(\log(\chi+1))$ and $O_d(N\log(\chi+1))$ clean auxiliary qudits.

The theorem is a nonuniform circuit-existence statement. Exact
parameters for the determinant maximizers and unitary extensions
are allowed as gate parameters; no efficient classical computation
of those parameters is asserted. The gate set consists of arbitrary
one- and two-qudit unitaries, as in the source, rather than a fixed
finite alphabet. The general exact circuit theorem and its
periodic-boundary specialization are formalized in
\nolinkurl{TNLean/MPS/MPU/ExactBoundedCutRankCircuit.lean}
(\nolinkurl{MPUCircuit.exists_clean_exact_circuit_of_bounded_cutCoefficientRank},
\nolinkurl{MPUCircuit.exists_clean_exact_circuit_of_periodic_mpo_unitary}).
The operator-norm approximation statements, with exact auxiliary erasure,
are formalized in
\nolinkurl{TNLean/MPS/MPU/ApproximateBoundedCutRankCircuit.lean}
(\nolinkurl{MPUCircuit.exists_clean_operatorNorm_approximate_circuit_of_bounded_cutCoefficientRank},
\nolinkurl{MPUCircuit.exists_clean_operatorNorm_approximate_circuit_of_periodic_mpo_unitary}).
These statements derive the complete circuit from the given finite
unitary and its consecutive cut-rank bound.

The exact theorem also proves the operator-norm approximation statement:
for every $\varepsilon>0$ the same circuit has zero error, with exactly
erased auxiliary outputs. This conclusion uses the arbitrary-gate model
of the theorem.

\subsection*{The transpose convention and the conditioning domain in the source}

Write the virtual matrices with the left bond as the row index and the
right bond as the column index, as in the tensor labels $m$ (left) and $n$
(right) of Styliaris--Trivedi--Cirac, arXiv:2508.08160v2, main-text
paragraph on MPUs with uniform bulk (source lines 695--722).
For an interval from site $j$ through site $k$, put
$B^{\mathbf a,\mathbf b}=A_j^{a_jb_j}\cdots A_k^{a_kb_k}$.
The two outward flag legs of the source's
\nolinkurl{eq:isometries_inhomo} have amplitudes
\[
 (V_{jk}^{[L_j,R_k]})_{(\mathbf a,\alpha,\beta),\mathbf b}
 =\sum_{u,v}(L_j)_{\alpha u}
       B^{\mathbf a,\mathbf b}_{uv}(R_k)_{\beta v}
 =(L_j B^{\mathbf a,\mathbf b}R_k^T)_{\alpha\beta}.
\]
The transpose is required because both flag operators act on outgoing
legs. The ordering of the two flag registers can be interchanged by a
fixed permutation; this does not change any norm or isometry assertion.

Let $P_t\succ0$ and $Q_t=r_t^{-1}P_t^{-1}$ be the compatible balanced
metrics at cut $t$ of Section~5. Its weighted interval amplitudes are
$\bigl(P_{j-1}^{1/2}B^{\mathbf a,\mathbf b}Q_k^{1/2}\bigr)_{\alpha\beta}$.
Consequently the source matrices for these same amplitudes are
\[
 L_j=P_{j-1}^{1/2},\qquad R_k=(Q_k^{1/2})^T.
\]
Transpose preserves positive definiteness. Both source matrices are
therefore positive and full rank. The simultaneous interval-isometry
identity proves their validity for every interval, including arbitrary
internally entangled physical inputs. At the endpoints,
$L_1=R_N=1$, so the full interval is exactly the original unitary.

At cut $k$, the source's joining operator is
$|00\rangle\langle I|(R_k^{-1}\otimes L_{k+1}^{-1})$; see the proof of
\nolinkurl{lem:merging}, source lines 2157--2182. Its coefficient
on the right flag $\beta$ of the left child and the left flag $\alpha$
of the right child is
\[
 \sum_t(R_k^{-1})_{t\beta}(L_{k+1}^{-1})_{t\alpha}
 =\bigl(R_k^{-T}L_{k+1}^{-1}\bigr)_{\beta\alpha}
 =\bigl(Q_k^{-1/2}P_k^{-1/2}\bigr)_{\beta\alpha}
 =\sqrt{r_k}\,\delta_{\alpha\beta}.
\]
Thus this operator is precisely the Bell reset used in Section~5,
with norm $r_k$. Equivalently, the source's squared norm becomes
\[
 \operatorname{Tr}\!\left[R_k^{-2}(L_{k+1}^{-2})^T\right]
 =\operatorname{Tr}\!\left[(Q_k^{-1})^T(P_k^{-1})^T\right]
 =\operatorname{Tr}(P_k^{-1}Q_k^{-1})=r_k^2.
\]
These equalities identify an explicit full-rank candidate; no
conditioning or trace identity is supplied as an extra assumption.

The domain of the infimum must be stated separately. The main-text
\nolinkurl{eq:def_q_k}, source lines 1375--1379, describes the
infimum over valid choices producing the displayed interval isometries.
For the unrestricted domain of all simultaneous valid positive metrics,
the construction above therefore gives $q_k^{\mathrm{iso}}\leq r_k$.
The appendix restatement, source lines 2123--2124, instead writes
$R_k^2\in\mathcal R_k$ and $L_{k+1}^2\in\mathcal L_{k+1}$ explicitly.
Those sets are defined from density operators on physical caps in
\nolinkurl{eq:def_LR}, source lines 1758--1816.
A positive matrix in the real affine hull of cap Grams need not itself
be a cap Gram, and the dual metric is not asserted to belong to a
physical cap set. Hence the argument does not bound that
cap-restricted infimum. The exact circuit theorem uses the explicit
Bell-reset norm $r_k$, independently of this difference between the
source's two stated optimization domains.

\subsection*{Physical cap conditioning can diverge for a single two-qubit gate}

Consider the two-qubit unitary
\[
 U_\theta=\operatorname{diag}(1,1,1,e^{i\theta}),\qquad 0<\theta<\pi.
\]
It is one allowed two-qudit gate in the circuit model of
Styliaris--Trivedi--Cirac, arXiv:2508.08160v2, opening of the
main-text paragraph on the main result (source line 820).
Nevertheless, the infimum restricted to physical cap density operators
in the appendix, source lines 2123--2124, diverges as
$\theta\downarrow0$. This is a calculation of an upper-bound parameter
for a particular construction; it is not a circuit lower bound.

Put $z=e^{i\theta}$ and use virtual row and column coordinates indexed by
$\{0,1\}$. A minimal open-boundary representation is given by
\[
\begin{array}{c|cccc}
 (a,b)&(0,0)&(0,1)&(1,0)&(1,1)\\ \hline
 A_1^{ab}&(1,0)&(0,0)&(0,0)&(0,1)\\[2pt]
 A_2^{ab}&\binom{1}{1}&\binom{0}{0}&\binom{0}{0}&\binom{1}{z}
\end{array}
\]
Thus $A_1^{ab}=\delta_{ab}e_b^T$ and
$A_2^{cd}=\delta_{cd}g_d$, where
$g_0=(1,1)^T$ and $g_1=(1,z)^T$. Their product is
$\delta_{ab}\delta_{cd}z^{bd}$, exactly the matrix entries of $U_\theta$.
The coefficient matrix
$\left(\begin{smallmatrix}1&1\\1&z\end{smallmatrix}\right)$
has determinant $z-1\ne0$, so the operator Schmidt rank is exactly two
throughout the stated nondegenerate interval.

The physical cap sets are defined in the source's
\nolinkurl{eq:def_LR}, source lines 1758--1816, by tracing the
physical outputs against arbitrary density operators on the cap inputs.
Here the off-diagonal entries of those density operators disappear,
because both local letters are diagonal on the physical indices.
Consequently the complete families, in the present virtual coordinates,
are
\[
 P_p=\begin{pmatrix}p&0\\0&1-p\end{pmatrix},\qquad
 Q_t=\begin{pmatrix}
 1&t+(1-t)\overline z\\t+(1-t)z&1
 \end{pmatrix},\qquad 0<p,t<1.
\]
Every displayed member is realized by a diagonal density operator, and
these are precisely the positive-definite members of the cap families.
In the outgoing-flag convention of the source,
$L_2^2=P_p$ and $R_1^2=Q_t^T$. In particular, the transpose appearing
in the source's joining norm gives
\[
 \operatorname{Tr}\!\left[R_1^{-2}(L_2^{-2})^T\right]
 =\operatorname{Tr}\!\left[(Q_t^{-1})^T(P_p^{-1})^T\right]
 =\operatorname{Tr}(P_p^{-1}Q_t^{-1}).
\]
There is no dimension normalization missing in these cap matrices:
$\operatorname{Tr}(P_pQ_t)=1$ for every $p,t$.

Since
\[
 \det Q_t=1-|t+(1-t)z|^2
          =4t(1-t)\sin^2(\theta/2),
\]
the complete inverse-trace expression is
\[
 \operatorname{Tr}(P_p^{-1}Q_t^{-1})
 =\frac{1}{4p(1-p)t(1-t)\sin^2(\theta/2)}.
\]
The inequalities $p(1-p)\le1/4$ and $t(1-t)\le1/4$, with equality
simultaneously at $p=t=1/2$, prove that the infimum is attained and equals
\[
 \boxed{\displaystyle q_{\mathrm{cap}}(U_\theta)
       =\frac{2}{\sin(\theta/2)}\sim\frac4\theta.}
\]
No limiting rank-deficient cap is needed for this minimum.

The same value can be checked in the canonical normalization of the
source, \nolinkurl{eq:choice_schmidt} and
\nolinkurl{eq:app:choice_schmidt}, source lines 1633--1640 and
1946--1952. The normalized Choi state has Schmidt values
\[
 c=\cos(\theta/4),\qquad s=\sin(\theta/4),\qquad c^2+s^2=1.
\]
Indeed, the two singular values of the coefficient matrix above are
$2c,2s$, and the normalized Choi state divides this matrix by two;
compare the normalization in source line 1393.
Let
\[
 H=\frac1{\sqrt2}\begin{pmatrix}
 e^{-i\theta/4}&e^{-i\theta/4}\\
 e^{i\theta/4}&-e^{i\theta/4}
 \end{pmatrix},\qquad X=\sqrt2H.
\]
Under the bond gauge $A_1\mapsto A_1X$ and
$A_2\mapsto X^{-1}A_2$, the cap matrices become
$\widetilde P=X^\dagger PX$ and
$\widetilde Q=X^{-1}QX^{-\dagger}$.
With $x=2p-1$ and $y=2t-1$, their complete families are
\[
 \widetilde P_x=\begin{pmatrix}1&x\\x&1\end{pmatrix},\qquad
 \widetilde Q_y=\begin{pmatrix}c^2&-i ycs\\i ycs&s^2\end{pmatrix},
 \qquad -1<x,y<1.
\]
Thus the maximally mixed caps give
$\widetilde L_2=I$ and $\widetilde R_1=\operatorname{diag}(c,s)$,
exactly the source's canonical scale. The inverse-trace becomes
$[(1-x^2)(1-y^2)c^2s^2]^{-1}$; its minimum is
$c^{-2}+s^{-2}=4/\sin^2(\theta/2)$.
More generally, the trace is unchanged by the gauge because
$\widetilde P^{-1}\widetilde Q^{-1}
 =X^{-1}P^{-1}Q^{-1}X$.

For comparison, the balanced metrics in the original coordinates are
$P_*=I/2$ and $Q_*=I$. The first local isometry has columns given by the
orthonormal flags $e_b$, while the second has columns $g_d/\sqrt2$ in
distinct physical output sectors; hence both are isometries and their
joining norm is exactly two. These metrics therefore give a valid
unrestricted candidate with $q_{\mathrm{iso}}\le2$.
They do not give a physical cap candidate: for $0<\theta<\pi$,
$t+(1-t)e^{-i\theta}$ cannot vanish, so $Q_*=I$ is outside the family
$\{Q_t\}$.
The distinction between the unrestricted valid-isometry domain in
source lines 1375--1379 and the physical-cap domain in lines
2123--2124 therefore has an observable mathematical effect, even for
this elementary rank-two unitary. A large physical-cap conditioning
number is compatible with an exact one-gate implementation.

\paragraph{The unrestricted two-site infimum.}
The balanced candidate need not minimize over all valid interval metrics.
In fact, this example also admits an explicit evaluation of that larger
optimization domain. With $Z=\operatorname{diag}(1,-1)$ and
$D_\theta=e^{-i\theta Z/4}$,
\[
 U_\theta=e^{i\theta/4}(D_\theta\otimes D_\theta)
       \bigl(cI\otimes I+i sZ\otimes Z\bigr).
\]
Local output unitaries and a global phase preserve all the Gram and
isometry conditions. Thus one may use the two letters $[I,Z]$ on the
first site and $[cI,i sZ]^T$ on the second. For this two-site chain the
only proper nonempty intervals are the two single sites; the full
interval is the original unitary, with trivial endpoint metrics.
The complete positive metrics making both single-site maps isometries
are
\[
 P=\begin{pmatrix}a&x\\x&b\end{pmatrix},\qquad
 Q=\begin{pmatrix}t&i y\\-i y&1-t\end{pmatrix},
\]
where $a,b>0$, $x,y\in\mathbb R$, $0<t<1$,
$x^2<ab$, $y^2<t(1-t)$, and $c^2a+s^2b=1$.
Indeed, the first site's rows $(1,1)$ and $(1,-1)$ force
$\operatorname{Tr}Q=1$ and $\operatorname{Re}Q_{01}=0$.
The second site's columns $(c,i s)^T$ and $(c,-i s)^T$ force
$c^2P_{00}+s^2P_{11}=1$ and $\operatorname{Im}P_{01}=0$.
There are no additional interval conditions in this two-site case.

The off-diagonal terms cancel in the inverse trace:
\[
 \operatorname{Tr}(P^{-1}Q^{-1})
 =\frac{b(1-t)+at}{(ab-x^2)[t(1-t)-y^2]}.
\]
For fixed $a,b,t$ this is minimized at $x=y=0$.
Setting $u=c^2a\in(0,1)$ then reduces the squared norm to
\[
 F(u,t)=\frac{c^2}{ut}
       +\frac{s^2}{(1-u)(1-t)}.
\]
The three-factor H\"older inequality gives
\[
\begin{split}
 c^{2/3}+s^{2/3}
 &=\left(\frac{c^2}{ut}\right)^{1/3}u^{1/3}t^{1/3}
   +\left(\frac{s^2}{(1-u)(1-t)}\right)^{1/3}
       (1-u)^{1/3}(1-t)^{1/3}\\
 &\le F(u,t)^{1/3}.
\end{split}
\]
Equality holds at
$u=t=c^{2/3}/(c^{2/3}+s^{2/3})$, strictly inside the permitted domain.
Consequently the unrestricted simultaneous-isometry infimum is
\[
 \boxed{\displaystyle q_{\mathrm{iso}}(U_\theta)
   =\bigl(c^{2/3}+s^{2/3}\bigr)^{3/2}.}
\]
It is strictly less than two for $0<\theta<\pi$ and tends to one as
$\theta\downarrow0$. This unrestricted optimization allows pairs for
which $\operatorname{Tr}(PQ)\ne1$; it is therefore distinct from the
positive-metric joining scheme normalized by
$\operatorname{Tr}(PQ)=1$, within which the balanced candidate attains
the optimal norm two. Neither calculation is an assertion about a
lower bound for arbitrary circuits.

\subsection*{Which polynomial circuit statement is established?}

Fix the physical dimension $d\ge2$. For $N,D\ge1$, let
$\mathcal U_d(N,D)$ denote the unitaries on $N$ $d$-dimensional sites
whose operator Schmidt ranks across all consecutive cuts are at most
$D$. No unitarity at other lengths and no lower bound on a nonzero
Schmidt value is imposed. Gates are arbitrary one- and two-qudit
unitaries, as in Styliaris--Trivedi--Cirac, arXiv:2508.08160v2, opening of the
main-text paragraph on the main result (source line 820).

For $A$ auxiliary qudits, write
$J_A\psi=\psi\otimes|0\rangle^{\otimes A}$.
An exact implementation is a unitary circuit $C$ satisfying
\[
 C J_A=J_AU.
\]
This is an equality on all initialized-auxiliary inputs, and remains an
equality after tensoring with the identity on any reference system.
It ensures exact auxiliary erasure. It does not require the stronger
ambient equality $C=U\otimes I$, which prescribes the action on
nonzero auxiliary inputs.

\paragraph{The exact theorem and its quantifiers.}
The determinant-normalization theorem of Section~5 establishes the
following bound. For each fixed $d$ there are constants $C_d,c_d>0$,
independent of $N,D$, the particular unitary, and its Schmidt values,
such that every $U\in\mathcal U_d(N,D)$ has an exact implementation with
\[
 G\le C_d(D+1)^{c_d}
             (N+1)^{c_d\log_2(D+1)},\qquad
 A\le C_d N\log_2(D+1).
\]
Increasing $c_d$ absorbs harmless numerical factors and the passage
to neighboring two-qudit gates on a line. Thus for every fixed $D$
the gate count is polynomial in $N$, with exactly clean auxiliaries.
The exponent may grow logarithmically with $D$.
The circuit and its exact gate parameters may depend on $U$:
this is a nonuniform existence assertion, not a claim that those
parameters can be computed efficiently from finite-precision tensor
entries. Uniformity of the numerical constants is distinct from an
algorithm for constructing the gates.

\paragraph{Joint dependence on length and bond dimension.}
The same estimate is a uniform quasi-polynomial bound in the numerical
parameters $N+D$. More precisely,
\[
 \log G\le O_d\!\left(
       \log(D+1)[1+\log(N+1)]\right)
 \le O_d\!\left(\log^2(N+D+2)\right),
\]
and hence
\[
 G\le \exp\!\left(O_d(\log^2(N+D+2))\right)
       =(N+D+2)^{O_d(\log(N+D+2))}.
\]
The constants in these expressions depend only on the fixed physical
dimension. This terminology refers to a gate-count bound in numerical
parameters; arbitrary complex tensor entries do not supply a finite
bit-length input model.
It does not establish a bound
\[
 G\le C_d(N+D+1)^{k_d}
\]
with an exponent $k_d$ independent of $D$.
For instance, putting $D=N$ into the established estimate gives
$\exp(O_d((\log N)^2))$, rather than a polynomial bound from this
argument. This observation is not a lower bound.

\textbf{Open stronger statement:} a bound polynomial jointly in $N$
and $D$ for exact arbitrary-gate implementations of every
$U\in\mathcal U_d(N,D)$, with exactly clean auxiliaries.
Retaining the present $O_d(N\log(D+1))$ auxiliary bound is a natural
version of this question. The argument here also does not prove the
version allowing an arbitrary polynomial number of auxiliaries.
``Open'' means that these stronger statements are not established by
the present proof or the cited local source; it is not a claim about
all possible circuit constructions.

\paragraph{Approximation and auxiliary erasure.}
For every $\varepsilon>0$, the exact circuit above already satisfies
\[
 \|C J_A-J_AU\|_{\mathrm{op}}=0\le\varepsilon,
\]
with the same gate and auxiliary counts and exact erasure.
Thus fixed-$D$ polynomial approximation in the arbitrary-gate model
is an immediate consequence; no additional dependence on
$\varepsilon$ is needed for this existence assertion.
The estimate is uniform over inputs, including inputs entangled with
an arbitrary reference system.

\textbf{Open stronger approximation statement:} for each fixed $d$,
constants $C_d,k_d$ independent of $N,D,\varepsilon,U$, giving
\[
 G\le C_d\bigl(N+D+\log(1/\varepsilon)+1\bigr)^{k_d},
 \qquad A\le C_dN\log_2(D+1),
\]
for every $U\in\mathcal U_d(N,D)$ and $0<\varepsilon\le1/2$, with
$\|C J_A-J_AU\|_{\mathrm{op}}\le\varepsilon$.
This weaker auxiliary requirement permits a residual nonzero
auxiliary component of norm at most $\varepsilon$.
The stronger approximation convention requires a unitary $V$ with
$C J_A=J_AV$ and $\|V-U\|_{\mathrm{op}}\le\varepsilon$, so that the
auxiliaries are erased exactly. A joint polynomial bound with this
stronger convention is also not proved here.
Polynomial dependence on $1/\varepsilon$ is another, weaker precision
requirement and must not be identified with polynomial dependence on
$\log(1/\varepsilon)$. The present quasi-polynomial exact bound by
itself proves neither joint polynomial formulation.

\paragraph{Changing the gate alphabet.}
The arbitrary-gate theorem permits exact continuous parameters.
A fixed finite gate alphabet cannot represent every such unitary
exactly: its finite circuits form a countable set.
Approximation with a chosen finite universal alphabet is a separate
formulation. If each allowed local gate can be approximated to error
$\delta$ with at most $L_d(\delta)$ alphabet gates, telescoping the
product gives an initialized-input error at most $\varepsilon$ by
choosing $\delta=\varepsilon/G$, with at most
$G L_d(\varepsilon/G)$ gates and the same auxiliary register allocation.
This implication is conditional on the stated local synthesis bound;
it does not assert an efficient algorithm for the determinant metrics.
The local approximation is in operator norm including scalar phase;
projective universality alone does not supply this hypothesis.
The conversion also does not by itself preserve exact auxiliary erasure.
In particular, such a gate-by-gate conversion does not remove the
existing logarithmic dependence of the exponent on $D$.

\paragraph{Small lengths and bond dimensions.}
The conventions $N,D\ge1$ and the shifts $N+1,D+1$ avoid spurious
zero logarithms. For $N=1$, one arbitrary one-qudit gate suffices;
if only two-qudit gates are counted, one clean padding qudit realizes
$U\otimes I$ as a single neighboring pair gate.
For $D=1$, every consecutive operator Schmidt rank is one, so the
unitary is a tensor product of on-site unitaries after distributing
its scalar phase. It therefore has an $O(N)$ implementation with no
auxiliaries in the one- and two-qudit model. The stated general
bounds include both cases, without asserting that their auxiliary
upper bounds are necessary.

\end{document}
